\documentclass[11pt,a4paper]{article}
\usepackage[utf8]{inputenc}
\usepackage[T1]{fontenc}
\usepackage[ngerman]{babel}
\usepackage{amsmath,amssymb}
\usepackage{graphicx}
\usepackage{booktabs}
\usepackage{xcolor}
\usepackage[margin=25mm]{geometry}
\usepackage{tcolorbox}
\usepackage{hyperref}

\definecolor{darkblue}{HTML}{16213e}
\definecolor{accentred}{HTML}{e94560}
\definecolor{keygreen}{HTML}{2e7d32}

\title{\textcolor{darkblue}{Dok. 0015\\Akustische Kopplung und Impedanzanpassung}}
\author{Harmonikabau — Technische Dokumentation}
\date{}

\begin{document}
\maketitle

\begin{tcolorbox}[colframe=red!60!black, colback=red!5!white]
\textbf{Hinweis:} Dieses Dokument bietet ein Erklärungsmodell.
Die angegebenen Zahlenwerte sind Größenordnungsabschätzungen ---
sie liefern keine konstruktiv verwertbaren Vorgaben.
Das Optimum wird gehört, nicht berechnet.
\end{tcolorbox}

\noindent\textit{Analyse der akustischen Kopplung zwischen Stimmzunge und Kammer als
Zwei-Port-System. Übergangsgeometrie, Verengung als akustische Masse,
Kopplungsschwerpunkt und Impedanzfehlanpassung.
Referenz: Dok.~0005--0013.}

\section{Die Kammer als Zwei-Port-System}

In den bisherigen Dokumenten (Dok.~0004--0008) wurde die Kammer als
Helmholtz-Resonator mit einer einzigen Öffnung (Klappe) behandelt.
Dok.~0013 hat gezeigt, dass die Zunge den Kanal nie vollständig verschließt ---
der Seitenspalt bleibt immer offen. Die Kammer hat daher \textbf{zwei Ports}:

\textbf{Port~1 (Zungenende):} Die Zunge pulsiert und erzeugt Druckimpulse.
Die effektive Fläche schwankt zwischen dem Seitenspalt
($A_\mathrm{min} \approx 10~\mathrm{mm}^2$) und der Schlitzfläche
($A_\mathrm{max} \approx 26~\mathrm{mm}^2$).

\textbf{Port~2 (Klappe):} Stationäre Öffnung mit $A_2 \approx 300~\mathrm{mm}^2$.
Hier strahlt der Schall ab.

Die Conductance (Leitwert) $G = A/l$ bestimmt, wie leicht Luft durch jeden
Port fließt. Port~1 (Zunge) hat zeitgemittelt etwa 70\,\% der Gesamt-Conductance ---
die Zunge ist der Hauptport, nicht die Klappe.

\begin{equation}
f_H = \frac{c}{2\pi}\sqrt{\frac{S_1/l_1 + S_2/l_2}{V}}
\end{equation}

\begin{tcolorbox}[colframe=keygreen, colback=green!5!white]
\textbf{Konsequenz:} Alle Berechnungen in Dok.~0004--0008, die $f_H$ nur aus der
Klappenöffnung ableiten, müssen als Untergrenze verstanden werden.
Die tatsächliche Kammerresonanz liegt höher.
\end{tcolorbox}

\section{Ist die Faltung akustisch sichtbar?}

Bei einem Bass mit 50\,Hz beträgt die Wellenlänge $\lambda = c/f = 6860$\,mm.
Die größte Kammerabmessung (Kammerlänge 120\,mm) ist nur 1,7\,\% von $\lambda$.
Eine Faltwand (30\,mm) ist 0,44\,\% von $\lambda$.

Wenn $d/\lambda < 0{,}1$, beugt sich die Welle vollständig um das Hindernis ---
es existiert akustisch nicht. Bei $d/\lambda > 0{,}5$ wirkt es als Reflektor.

Bei 50\,Hz liegen \textbf{alle} Kammergeometrien tief im Bereich $d \ll \lambda$.
Erst ab $\approx 1400$\,Hz (28.~Oberton) beginnt die Faltwand die Grenze
$d/\lambda = 0{,}1$ zu erreichen.

Die Analogie zum Lichtwellenleiter bestätigt dies: Ein glattes, rundes Rohr
mit konstantem Querschnitt hat keinen Reflexionsfaktor --- die Welle folgt
der Kurve, solange der Biegeradius groß gegen die Wellenlänge ist.

\textbf{Coltman-Korrektur:} Coltman kompensiert nicht die Biegung selbst,
sondern die Querschnittsänderung und Pfadlängendifferenz an der Faltstelle.
$Z = \rho c / A$ bleibt konstant bei konstantem Querschnitt.
Die Pfadlängenkorrektur skaliert mit $d^2/r$ --- bei $d/\lambda < 0{,}01$
ist sie vernachlässigbar.

\begin{tcolorbox}[colframe=keygreen, colback=green!5!white]
\textbf{Fazit:} In der Basskammer spielt die Faltung für den Grundton und die
ersten Obertöne keine Rolle. Die Anpassung geschieht an den Öffnungen
(Port~1 und Port~2), nicht an der Faltung.
\end{tcolorbox}

\section{Verengung = akustische Masse}

Jede Querschnittsverkleinerung im Kanal wirkt als konzentrierte akustische Masse:
$m = \rho \times l / A$. Das gilt unabhängig davon, ob die Verengung in einem
geraden Rohr, an einer Faltstelle oder am Übergang Schlitz$\to$Kammer sitzt.

Die Coltman-Kompensation bei gefalteten Orgelpfeifen ist im Kern eine
Verengung an der Umlenkstelle. Der Mechanismus ändert sich
(Wellenleiter vs.\ Lumped Element), das physikalische Element
(Verengung = Masse) bleibt dasselbe.

\section{Druckaufbau vor der Engstelle}

Der Druckimpuls der Zunge staut sich vor der Engstelle am Übergang
Schlitz$\to$Kammer. Drei Szenarien:

\textbf{Zu weit:} Druck entweicht sofort ins Volumen --- schwache Rückwirkung,
geringe Kopplung.

\textbf{Optimal:} Druck staut sich kurz auf --- maximale Rückwirkung,
beste Kopplung.

\textbf{Zu eng:} Impuls wird überwiegend reflektiert --- Energie kommt
nicht ins Volumen.

\section{Impedanzanpassung: Übergangsgeometrie als Transformator}

Die Beobachtung, dass verschiedene Übergangsgeometrien die \textbf{Ansprache}
ändern, aber nicht die \textbf{Filterwirkung}, lässt sich sauber als
Anpassungsproblem modellieren:

\textbf{Quelle:} Der Balgdruck $p_\mathrm{Balg}$ mit einem Innenwiderstand $R_i$,
der vom Strömungswiderstand des Systems abhängt. $R_i \propto 1/A^2$ (Bernoulli).

\textbf{Last:} Die Zunge mit mechanischer Impedanz $R_L$ (fix, von Steifigkeit
und Masse bestimmt, Dok.~0012).

\textbf{Transformator:} Die Übergangsgeometrie am Zungenende bestimmt,
wie effizient die Energie der Strömung in Zungenenergie umgesetzt wird.

\begin{equation}
P / P_\mathrm{max} = \frac{4 \, R_L \, R_i}{(R_i + R_L)^2}
\end{equation}

Maximum bei $R_i = R_L$ --- klassische Anpassungsbedingung.
Zu eng: $R_i \gg R_L \to$ Blockade. Zu weit: $R_i \ll R_L \to$ Kurzschluss.
Optimal: $R_i = R_L \to$ maximaler Leistungstransfer Balg$\to$Zunge.

\begin{tcolorbox}[colframe=keygreen, colback=green!5!white]
\textbf{Entscheidend:} Die Filterwirkung ($f_H$, Kerbfilter) hängt vom
Resonator ab --- Kammervolumen und Öffnungsfläche. Der Transformator
(Übergangsgeometrie) ändert die Anpassungseffizienz, nicht die Resonanz.
Das ist genau das, was man beobachtet: andere Ansprache, gleiche Klangfarbe.
\end{tcolorbox}

\section{Ersatzschaltbild: Zwei gekoppelte Schwingkreise}

Die Trennung von Ansprache und Filterwirkung lässt sich als
elektroakustische Analogie exakt darstellen. Zunge und Kammer sind je
ein Serienschwingkreis (RLC), verbunden durch ein Kopplungselement:

\textbf{Kreis 1 (Zunge):} Masse $m \to$ Induktivität $L_1$.
Steifigkeit $k \to$ Kapazität $C_1 = 1/k$.
Dämpfung $\to$ Widerstand $R_1$.
$f_1 = 1/(2\pi\sqrt{L_1 C_1}) \approx 50$~Hz.

\textbf{Kreis 2 (Kammer):} Akustische Masse $\rho l/A \to L_2$.
Volumennachgiebigkeit $V/(\rho c^2) \to C_2$.
$f_H = 1/(2\pi\sqrt{L_2 C_2})$.

\textbf{Kopplung:} Gegenseitige Impedanz $M = k \times \sqrt{L_1 L_2}$.
Der Kopplungskoeffizient $k$ (0--1) wird von der Übergangsgeometrie bestimmt.

Die Transferfunktion $H(f) = I_2 / V_1$ zeigt zwei Resonanzpeaks (bei $f_1$ und $f_H$).
Entscheidend: $k$ ändert die \textbf{Höhe} der Peaks (Energietransfer = Ansprache),
aber nicht ihre \textbf{Position} (Frequenz = Filterwirkung). Umgekehrt
verschiebt $C_2$ (Kammervolumen) den $f_H$-Peak, ohne die Ansprache zu ändern.

\begin{tcolorbox}[colframe=keygreen, colback=green!5!white]
\textbf{Berechenbar:} Mit den physikalischen Werten (Masse, Steifigkeit,
Volumen, Öffnungsfläche) ergibt sich die Transferfunktion quantitativ.
$k$ bestimmt Ansprache, $C_2$ bestimmt Filterwirkung ---
diese Parameter sind unabhängig.
\end{tcolorbox}

\section{Klappenposition --- Feinkorrektur}

Die Klappenöffnung kann in der Praxis um einige Millimeter bis maximal
einen Zentimeter vom Kammerende versetzt werden. Die Verschiebung ändert
die effektive akustische Länge des Halses und damit $f_H$.

Im Praxisbereich (0--15\,mm) verschiebt sich $f_H$ um wenige Prozent.
Die Klappenposition ist eine Feinkorrektur --- der Haupthebel liegt am
Zungenende (Port~1), wo 70\,\% der Conductance sitzen.

\section{Grenzen des Modells}

Das Zwei-Port-Modell zeigt Zusammenhänge und Richtungen.
Es sagt nicht, wo das Optimum liegt.

\textbf{(a) Nichtlinearität:} Die Zunge schaltet periodisch zwischen
kleiner und großer Öffnung. Die tatsächliche Impedanzmodulation
(Dok.~0013) ist hochgradig nichtlinear.

\textbf{(b) Klangziel:} Das ``optimale'' Spektrum hängt vom gewünschten
Klangcharakter ab. Die Berechnung zeigt die Landkarte, den Weg geht das Ohr.

\textbf{(c) Wechselwirkungen:} Spaltweite (Dok.~0010), Steifigkeit (Dok.~0012),
Kammervolumen, Klappenöffnung und Übergangsgeometrie wirken zusammen.

\begin{tcolorbox}[colframe=red!60!black, colback=red!5!white]
\textbf{Dieses Dokument ist ein Erklärungsmodell, keine Konstruktionsanleitung.
Konstruktionsentscheidungen erfordern nach wie vor praktische Tests und Erfahrung.}
\end{tcolorbox}

\section{Wo koppelt die Zunge? --- Schwerpunkt des Volumenstroms}

Die Annahme, dass die Kopplung am absoluten Ende der Zunge (Spitze) erfolgt,
ist eine Vereinfachung. Die Zunge ist ein Cantilever --- ihre Schwingform
folgt der ersten Euler-Bernoulli-Biegemode.

Der Volumenstrom pro Längeneinheit $\mathrm{d}Q/\mathrm{d}x \propto y(x)$
steigt steil zum Ende hin an. Der gewichtete \textbf{Schwerpunkt des
Volumenstroms} liegt bei $x/L = 0{,}73$ --- nicht an der Spitze ($x/L = 1{,}0$).
50\,\% des gesamten Volumenstroms entstehen erst in den letzten 23\,\%
der Zungenlänge.

Für eine Bass-Zunge mit $L = 70$\,mm liegt der Schwerpunkt bei 51\,mm
von der Einspannung --- ca.\ 19\,mm vor der Spitze.

\subsection{Impedanzfehlanpassung}

Die mechanische Impedanz der Zunge und die akustische Impedanz der Luft
unterscheiden sich um Größenordnungen. Die Zunge ist ein
\textbf{Strom-Generator} (hohe Impedanz) --- sie erzwingt einen Volumenstrom,
den die Luft aufnimmt.

$Z_\mathrm{Zunge}$ liegt Faktor 50--500 über $Z_\mathrm{Luft}$.
Die Kopplung ist \textbf{nicht angepasst} im Sinne eines maximalen
Leistungstransfers ($Z_\mathrm{Quelle} = Z_\mathrm{Last}$).
Das ist physikalisch richtig so: Die Zunge ist selbsterregt und schwingt
auf ihrer Eigenfrequenz, nicht auf der Kammerresonanz. Die Luft wird
``mitgenommen'', nicht ``angetrieben''.

Dies steht im Einklang mit Dok.~0010: Über 97\,\% der Energieabgabe erfolgt
über die Luftkopplung, weniger als 3\,\% über die mechanische Einspannung.

\subsection{Konsequenz: Position der Plattenkante}

Wenn die Stimmplatte die volle Zungenlänge abdeckt, koppelt der gesamte
Schlitz in die Kammer. Ragt die Zunge über die Plattenkante hinaus (Aufbiegung),
geht der energiereichste Teil --- die letzten 20--30\,\% der Zunge ---
an der Kammer vorbei, direkt in den freien Raum.

Die Position der Plattenkante relativ zur Zunge bestimmt daher, welcher Anteil
des Volumenstroms in die Kammer eingekoppelt wird und welcher Anteil
direkt abstrahlt.

\section{Zusammenfassung}

\begin{enumerate}
\item \textbf{Zwei-Port-System:} Die Kammer ist beidseitig offen.
  Port~1 (Zunge) trägt $\approx 70\,\%$ der Conductance.
  Das Ein-Port-Modell unterschätzt $f_H$.
\item \textbf{Faltung unsichtbar:} Bei 50\,Hz ist die größte Kammerabmessung
  1,7\,\% von $\lambda$. Die Kammerform spielt keine Rolle ---
  nur das Volumen zählt.
\item \textbf{Impedanzanpassung:} Die Übergangsgeometrie wirkt als Transformator
  zwischen Balgdruck (Quelle) und Zunge (Last). Maximaler Leistungstransfer
  bei $R_i = R_L$. Die Filterwirkung ($f_H$) hängt vom Resonator ab,
  nicht vom Transformator.
\item \textbf{Ersatzschaltbild:} Zunge und Kammer bilden zwei gekoppelte
  RLC-Schwingkreise. $k$ bestimmt die Ansprache (Energietransfer),
  $C_2 = V/(\rho c^2)$ bestimmt die Filterwirkung ($f_H$).
  Diese Parameter sind unabhängig.
\item \textbf{Klappenposition = Feinkorrektur:} Wenige Prozent $f_H$-Verschiebung.
  Der Haupthebel liegt am Zungenende.
\item \textbf{Schwerpunkt bei $x/L = 0{,}73$:} Die Kopplung sitzt nicht an
  der Zungenspitze, sondern $\approx 27\,\%$ davor.
  Die Position der Plattenkante bestimmt die Einkopplung.
\item \textbf{$Z_\mathrm{Zunge} \gg Z_\mathrm{Luft}$:} Die Zunge ist ein
  Strom-Generator. Die Impedanzen sind nicht angepasst --- die Luft wird
  mitgenommen, nicht angetrieben.
\item \textbf{Optimum wird gehört:} Die Berechnung zeigt Zusammenhänge und
  Richtungen. Erfahrung und Gehörschulung bleiben der Schlüssel zum Erfolg.
\end{enumerate}

\bigskip
\textit{Die Zunge speist, die Kammer formt, die Klappe strahlt ab.
Die Geometrie am Übergang bestimmt die Kopplung.}

\end{document}
